\documentclass[11pt,a4paper]{article}

% ======================================================================
%  The Maintenance Constraint: How Resource Boundaries Shape Cognitive
%  Availability. Version 2 (July 2026). v1: December 2025
%  (ai.viXra:2512.0084v1).
%  This is the outermost, philosophy-of-mind paper of the coherence-
%  maintenance series: no theorems of its own; its single load-bearing
%  technical input is the maintenance inequality, imported as an anchor.
%  Main changes (Appendix A):
%   1. The anchor inequality P_extra >= k_B T Cdot_loss is cited in its
%      CORRECTED form: v1 cited it as an unqualified extra-power bound, but
%      the extra-power statement is meaningful only against thermodynamically
%      efficient baselines, or unconditionally as a total-power / entropy-
%      production floor; the anchor references are updated to the v2 series
%      (2512.0081 v2 and, through it, 2512.0061 v2). The philosophical
%      argument is unchanged, and is if anything cleaner: the passive-
%      stability special case is exactly the free-baseline case.
%   2. Full series references added (v1 cited only [1], an unnumbered
%      preprint).
%   3. Non-claims retained verbatim; scope tightened where the anchor is
%      invoked.
% ======================================================================

\usepackage[utf8]{inputenc}
\usepackage[T1]{fontenc}
\usepackage{amsmath,amssymb}
\usepackage[margin=1.05in]{geometry}
\usepackage[colorlinks=true,linkcolor=blue,citecolor=blue,urlcolor=blue]{hyperref}

\newcommand{\kB}{k_{\mathrm{B}}}
\newcommand{\Pextra}{P_{\mathrm{extra}}}
\newcommand{\Pmin}{P_{\min}}
\newcommand{\Pavail}{P_{\mathrm{avail}}}
\newcommand{\Closs}{\dot{C}_{\mathrm{loss}}}

\title{The Maintenance Constraint:\\
How Resource Boundaries Shape Cognitive Availability}
\author{Lluis Eriksson\\ \small Independent Researcher\\
\small \texttt{lluiseriksson@gmail.com}}
\date{July 4, 2026 \quad (v2; v1: December 2025) \quad ai.viXra:2512.0084}

\begin{document}

\maketitle

\begin{center}
\fbox{\parbox{0.94\textwidth}{\small
\textbf{Status note.} This is the outermost, philosophy-of-mind paper of the
coherence-maintenance series (ai.viXra:2512.0060/0061/0064/0070/0071/0072/
0073/0081, all at v2, 2026). It contains no theorems of its own; its single
load-bearing technical input is the maintenance inequality, imported as an
\emph{anchor}. \textbf{v2}: that inequality is cited in its corrected form
--- the incremental (extra-power) statement holds against thermodynamically
efficient baselines, with an unconditional total-power / entropy-production
floor underneath it, per the corrected companion papers \cite{Anchor,Work}
(v2); v1 cited the unqualified extra-power form, whose companion-v1
formulation was later found vacuous. The philosophical argument is
unaffected. \textbf{Non-claims:} no collapse mechanism, no Born-rule
derivation, no claim about phenomenological consciousness or quantum
coherence in brains.}}
\end{center}

\begin{abstract}
Cognitive systems are resource-limited, but ``resource limitation'' is often
invoked without distinguishing one-shot costs (forming a representation) from
sustained costs (keeping it usable under noise). We argue that the
availability of internal state features for control, integration, and report
is constrained by their \emph{maintainability} under finite budgets. As a
technical anchor we cite a companion preprint deriving an operational
maintenance inequality in explicit thermodynamic control models: incremental
maintenance power has a non-arbitrary lower bound tied to dynamical fragility
(in its corrected form, against efficient baselines, with an unconditional
entropy-production floor beneath). This motivates an operational cut: a
feasibility boundary separating maintainable from unmaintainable state
features. We develop an auditable bridge argument (maintainability $\to$
stability $\to$ availability), propose a neutrality-friendly principle of
\emph{maintenance-feasibility bias}, and show how it can be incorporated into
active inference (the Free-Energy Principle) as a maintenance penalty or
constraint, aligning secondarily with Global Workspace accounts of access
stability. We address objections and offer falsifiable predictions for
synthetic agents and neuromorphic systems, plus an explicitly exploratory
psychophysics subsection framed in terms of reportability and stability
rather than phenomenology. We do not propose collapse mechanisms, do not
derive the Born rule, and make no claims about phenomenological
consciousness.
\end{abstract}

\section{Introduction}

Many debates in philosophy of mind and cognitive neuroscience presuppose that
internal variables can be treated as stable quantities available for
downstream computation, control, and report. Yet biological and artificial
agents operate under noise, drift, coupling, and finite budgets (metabolic
power, bandwidth, control precision, compute). This raises a neglected
question: which internal state features are feasible to keep \emph{available}
over task-relevant timescales?

A common move is to invoke ``resource limits'' in broad terms. While true,
this is too coarse to constrain architectures. Our aim is to sharpen it into
an operational constraint: \emph{availability is limited by maintainability}.
Some state features may be representable in principle but too fragile or too
costly to keep stably available under realistic budgets; this shapes which
representations dominate behavior and report.

\paragraph{States vs.\ state features.} We speak of \emph{state features}
rather than exact microstates: what must be stabilized is typically a
task-relevant feature (a coarse-graining, a decoded variable, a functional
manifold, a correlation pattern), not an entire microscopic configuration.

Two methodological commitments guide the paper. (1) We treat ``availability''
as a functional notion (access for downstream use), not a claim about
phenomenology; this targets access-level cognition and report, not the hard
problem \cite{Block95,Chalmers95,Chalmers96}. (2) We keep the physics anchor
separate from philosophy-of-mind claims: the anchor supplies an example of a
nontrivial lower bound on maintenance cost in an explicit model; it does not
imply that brains implement quantum coherence, nor that consciousness depends
on quantum effects.

\paragraph{Non-claims (explicit).} We do not propose collapse mechanisms, do
not derive the Born rule, and make no claims about phenomenological
consciousness. The term ``cut'' is used operationally --- a feasibility
boundary under resource constraints, not an ontological boundary in nature.

\section{Existing approaches to physical constraints on cognition}

That cognition is physically constrained is not new; what is often missing is
a principled separation between constraints on \emph{computation/
representation} and constraints on \emph{maintenance} of state features over
time under noise and coupling. Computational and informational limits (bounded
rationality \cite{Simon55}, resource-rational analysis \cite{LiederGriffiths20},
Marr's levels \cite{Marr82}) bound what can be \emph{computed} or
\emph{encoded}. Metabolic-cost accounts \cite{Laughlin98,AttwellLaughlin01}
bound the \emph{energy} of signaling. The present proposal is orthogonal to
both: it concerns the ongoing cost of \emph{holding} a state feature usable,
which neither a computational-complexity bound nor a per-spike energy budget
directly captures.

\section{Technical anchor: a maintenance inequality and an operational cut}

We summarize the anchor at a level sufficient for the philosophical argument,
without requiring the full physics proof \cite{Anchor,Work}.

\subsection{The anchor inequality (corrected form)}

In an explicit thermodynamic control model (finite-dimensional system,
Markovian uncontrolled dynamics, a work reservoir/battery interacting via
energy-conserving global unitaries with a thermal environment at temperature
$T$), one defines an incremental maintenance power $\Pextra(\rho)$ isolating
the steady cost of stabilizing certain state features at fixed populations.
The anchor result is the explicitly qualified pair of lower bounds
\begin{equation}\label{eq:anchor}
\Pmin(\rho) \;\ge\; \kB T\,\sigma(\rho),
\qquad\qquad
P^{\mathrm{eff}}_{\mathrm{inc}}(\rho) \;\ge\; \kB T\,\Closs(\rho)
\quad\text{in the efficient/free-baseline regimes},
\end{equation}
where $C(\rho)$ is a relative-entropy-based coherence functional,
$\Closs(\rho)$ its instantaneous loss rate under the uncontrolled dynamics,
and $\sigma(\rho) \ge 0$ the target's entropy production rate. Here
$P^{\mathrm{eff}}_{\mathrm{inc}}$ denotes the \emph{corrected} incremental
quantity: extra power relative to thermodynamically efficient
population-maintenance baselines, or the free-baseline case where
$\Pextra = \Pmin$. No unqualified paired-infimum statement is intended.

\paragraph{Provenance of \eqref{eq:anchor} (v2).} Both bounds are proved in
the companion papers \cite{Work,Anchor} (v2): the first is unconditional; the
second holds against thermodynamically efficient population-maintenance
baselines and becomes assumption-free when the stabilized feature's
``populations'' are already stationary (the passive-stability case of
\S4.5). An earlier companion formulation defined the incremental cost by an
unconstrained comparison and was vacuous; the corrected form is the one used
here. The
\emph{philosophically important} feature is structural and survives the
correction intact: maintenance has a non-arbitrary lower bound tied to
dynamical fragility.

We emphasize: the anchor is a \emph{benchmark case study} of maintainability
costs, not a claim about neural implementation. The argument does not require
that the brain maintains quantum coherence; the relevant import is the
existence of principled maintenance lower bounds in explicit physical models.

\subsection{The operational cut}

Given a finite budget $\Pavail$, define a feasibility boundary:
\begin{quote}
\emph{Operational cut:} the boundary between state features whose relevant
maintenance power $P_{\mathrm{maint}}$ --- total ($\Pmin$), incremental
against an efficient baseline ($P^{\mathrm{eff}}_{\mathrm{inc}}$), or
free-baseline, depending on the control task --- lies below the available
budget $\Pavail$, and those for which it does not.
\end{quote}
Using the task-appropriate $P_{\mathrm{maint}}$ rather than a bare
$\Pextra$ avoids transporting the incremental notion into cognitive contexts
where no clean ``populations baseline'' exists.
This cut is operational and agent-relative --- it depends on available
resources and the uncontrolled dynamics --- not an ontological division in
nature. Its cognitive role emerges once maintainability is connected to
stability and availability.

\subsection{Methodological status}

Our use of the anchor is methodological, not neurobiological. It provides an
existence proof that, in explicit control models, holding certain features
against uncontrolled dynamics entails non-arbitrary lower bounds tied to
fragility. We do not infer that brains instantiate the anchor's microphysics;
we use it to motivate a general stance: in any noisy substrate, availability
of fine-grained internal variables over task-relevant horizons may require
ongoing stabilization whose cost is constrained by the substrate's loss rates
and coupling structure. The claim defended here concerns the
architecture-shaping role of such feasibility boundaries, not the
universality of any specific inequality.

\section{From maintainability to cognitive availability}

The central bridge is stated as an explicit, auditable argument. The claim is
\emph{not} that maintainability is consciousness; rather, maintainability
constrains which internal state features can stably participate in cognition
in the sense of being available for control, integration, and report.

\subsection{Three notions}

\textbf{Maintainability (control-feasibility).} A target feature is
maintainable, relative to the uncontrolled dynamics and control primitives, if
there exists an admissible strategy keeping the system in (or near) a
functionally acceptable set for the requisite duration within budget. In the
anchor model, maintainability is constrained by a lower bound on incremental
maintenance power.

\textbf{Stability (dynamical robustness).} A feature is stable over a horizon
if, under actual perturbations and couplings, it does not drift rapidly in
ways that destroy its functional role. Stability concerns trajectory
timescales, not instantaneous existence of an encoding.

\textbf{Availability (cognitive-functional).} A feature is available if it can
be reliably used for downstream tasks: guiding action, updating beliefs,
broadcasting to other subsystems, supporting report. Following the
access/phenomenal distinction \cite{Block95,DehaeneNaccache01}, we treat
availability as separable from phenomenology; this paper concerns availability
and control, not qualia.

\textbf{Maintenance resources (non-quantum reading).} Measurable ongoing costs
--- metabolic power, control effort, communication bandwidth,
precision/feedback gain --- required to keep relevant variables within
functional bounds.

\subsection{Operational proxies}

\textbf{Stability proxies.} Over a horizon $T$: (i) mean deviation from a
functional set $K$, (ii) dwell-time in $K$ before exit, (iii) decay time of an
autocorrelation of a task-relevant variable.

\textbf{Availability proxies.} Downstream readout at latency $\tau$: (i)
decoding accuracy under noise, (ii) control-performance degradation when the
variable is perturbed, (iii) mutual information between the variable and
downstream actions within the relevant window. The maintenance constraint
predicts systematic trade-offs between these proxies and measured maintenance
expenditure under budget manipulations.

\subsection{The dependency chain}

\textbf{P1 (Availability requires task-relevant persistence).} To be available
for control, integration, or report, a feature must persist on timescales
comparable to downstream latencies; a microstate overwritten before it can
influence downstream processes is not available in the relevant sense.

\textbf{P2 (Persistence under noise often requires maintenance).} Uncontrolled
dynamics degrade fine-grained features toward default configurations; unless
the target lies in a naturally stable manifold or is architecturally protected,
persistence typically requires maintenance (feedback stabilization, metabolic
expenditure, error correction, active control).

\textbf{P3 (Maintenance can have unavoidable lower bounds).} Under explicit
thermodynamic bookkeeping, sustaining deviations from default relaxation can
impose non-zero lower bounds on ongoing resource use. The anchor provides a
paradigmatic case (quantum coherence) as an operational benchmark: in its
corrected reading,
$P^{\mathrm{eff}}_{\mathrm{inc}}(\rho) \ge \kB T\,\Closs(\rho)$ against
efficient baselines (unconditional total-power floor beneath), linking
incremental maintenance power to intrinsic dynamical fragility.

\textbf{C (Availability is resource-constrained).} From P1--P3: the set of
features available to an agent is constrained by its resource budget; some
candidates exist as encodings in principle but are unavailable in practice
because too costly to maintain under actual noise and coupling.

\begin{quote}
\textbf{Maintenance-feasibility bias:} when multiple internal encodings serve
comparable functional roles, finite-budget agents tend to adopt encodings
whose stability is achievable within available maintenance resources.
\end{quote}

\emph{What the conclusion does not say.} Not that cognition is nothing but
maintenance cost, nor that phenomenology is determined by the bound: it is a
constraint on availability and stability of internal variables under finite
budgets.

\subsection{Loading vs.\ holding: work vs.\ power}

A central operational distinction separates \emph{loading} a representation
(one-shot shaping cost to enter a target state) from \emph{holding} it
(sustained cost to keep it from drifting). This disentangles ``the system can
encode $X$'' from ``the system can stably use $X$''. An encoding may be cheap
to load but expensive to hold, or vice versa, predicting dissociations:
transient high-fidelity microstates unavailable for report, versus
lower-fidelity states that dominate behavior because they persist. (In the
anchor, this is exactly the one-shot-work vs maintenance-power distinction
made precise in \cite{Work,HCut} v2.)

\subsection{Passive stability as a special case}

Not all stable representations require active maintenance: strong attractors,
passive memory traces, and structural constraints yield persistence without
ongoing work. This is a special case where the uncontrolled dynamics already
keep the state near target, so incremental maintenance cost is negligible ---
in the anchor model, a vanishing loss rate for the relevant feature, making
the bound trivial. Notably, this is \emph{precisely} the free-baseline case
in which the corrected inequality \eqref{eq:anchor} holds without further
assumptions \cite{Work} (v2): the correction and the passive-stability limit
are the same statement seen from two sides. The maintenance-feasibility bias
reduces to the correct special case: availability is unconstrained by
maintenance when dynamics are already aligned.

\section{Connections to cognitive frameworks}

\subsection{Active inference / FEP with a maintenance penalty}

Active inference proposes that agents minimize variational free energy as a
proxy for surprise \cite{Friston10,ParrFriston19}. The maintenance constraint
adds an implementational dimension: even if a representational state reduces
variational objectives, it may be too fragile to hold under realistic
perturbations. This motivates a multi-objective trade-off
\begin{equation}
\min_{q,\pi}\; \mathbb{E}_\pi\!\big[F_{\mathrm{var}}(q)\big]
+ \lambda\, P_{\mathrm{maint}}(q,\pi), \qquad \lambda > 0,
\end{equation}
with $P_{\mathrm{maint}}$ the task-appropriate maintenance power of \S3.2
(in physical instances, the corrected efficient-/free-baseline
$P^{\mathrm{eff}}_{\mathrm{inc}}$): fragile
encodings that lower $F_{\mathrm{var}}$ either raise measured maintenance
power or are abandoned once the budget is capped.

\subsection{Global Workspace: access stability}

Global Workspace accounts \cite{Baars88,DehaeneNaccache01,Dehaene14} describe
broadcast/access dynamics for reportable content. The maintenance constraint
is a secondary alignment: broadcast requires that content remain stable long
enough to be accessed, so access stability is a maintainability condition on
workspace occupancy, not a claim about the workspace's phenomenal status.

\section{Objections}

\textbf{``If the brain is classical, the quantum anchor is irrelevant.''} The
anchor is not a claim about neural microphysics; it demonstrates in an
explicit model that (a) maintenance can carry principled lower bounds, and (b)
those bounds are a matter of physics, not engineering inefficiency. The
philosophical constraint is substrate-general: any noisy substrate with
nonzero loss rates for the relevant feature inherits a maintainability
condition, whatever its microphysics.

\textbf{``This is just metabolic cost renamed.''} No: metabolic-cost accounts
bound signaling energy; the maintenance constraint bounds the ongoing cost of
\emph{holding a feature usable}, which can be nonzero even when signaling is
cheap, and zero (passive stability) even when signaling is expensive.

\textbf{``Maintainability is not availability.''} Correct --- it is a
\emph{necessary} condition via the P1--P3 chain, not identical to it; the
paper claims constraint, not identity.

\section{Falsifiable implications}

\textbf{Synthetic/neuromorphic agents.} Under explicit energy/compute budgets,
policies achieving lower $F_{\mathrm{var}}$ via fragile encodings will either
increase measured maintenance power or revert to robust/low-cost encodings
once available power is capped --- a directly testable budget-manipulation
prediction. \textbf{Dissociation prediction.} Cheap-to-load,
expensive-to-hold features should be underrepresented in report relative to
their encodability, and the gap should widen as budgets tighten.

\paragraph{Exploratory psychophysics (framed as reportability/stability, not
phenomenology).} If report requires access stability, manipulations that raise
the maintenance cost of a decoded feature (increased noise/coupling on its
substrate) should reduce its reportability at fixed encodability, with the
effect scaling with the feature's intrinsic loss rate. This subsection is
explicitly exploratory and makes no phenomenological claim.

\section{Conclusion}

``Resource limitation'' in cognition is sharpened here into a specific,
operational constraint: availability is limited by maintainability, and
maintainability can carry principled lower bounds under finite budgets. A
companion technical preprint provides a benchmark instance in which
maintenance has a non-arbitrary lower bound tied to dynamical fragility
\cite{Anchor,Work} (v2). The resulting maintenance-feasibility bias is a
neutrality-friendly, falsifiable organizing principle for which internal
representations finite-budget agents can keep available --- with the physics
anchor kept firmly separate from, and not load-bearing for, any claim about
phenomenology.

\appendix

\section{Changes relative to version 1}

\begin{enumerate}
\item \textbf{Anchor corrected, argument preserved.} v1 cited the maintenance
inequality \eqref{eq:anchor} as an unqualified extra-power bound (its
reference [1], an unnumbered companion preprint). The companion's v1
extra-power definition was subsequently found vacuous (an infimum over
unlinked strategy pairs equals $-\infty$); the corrected hierarchy
\cite{Work,Anchor} (v2) reads \eqref{eq:anchor} as an efficient-baseline
incremental bound with an unconditional entropy-production floor beneath.
This paper cites the corrected form. The philosophical bridge (P1--P3,
maintenance-feasibility bias) is unchanged and, if anything, sharper: the
passive-stability special case (\S4.5) is now identified with the corrected
inequality's free-baseline regime.
\item \textbf{Series references added.} v1 cited only the single anchor
preprint; v2 cites the full v2 series and situates the anchor within it.
\item \textbf{Non-claims retained verbatim}; scope tightened where the anchor
is invoked (\S3.1, P3).
\end{enumerate}

\begin{thebibliography}{99}\small

\bibitem{Anchor} L.~Eriksson, \emph{Operational Coherence Maintenance and the
Quantum--Classical Boundary: Formal Definitions, Falsifiable Protocols, and
an Outlook for Cognitive Systems}, ai.viXra:2512.0081, v2 (2026; v1: 2025).
Companion repository, \texttt{papers/0081-quantum-classical-boundary}.

\bibitem{Work} L.~Eriksson, \emph{The Conditional Maintenance Work Theorem},
ai.viXra:2512.0061, v2 (2026; v1: 2025). Companion repository,
\texttt{papers/0061-maintenance-work}.

\bibitem{HCut} L.~Eriksson, \emph{The Heisenberg Cut as a Resource Boundary},
ai.viXra:2512.0064, v2 (2026; v1: 2025). Companion repository,
\texttt{papers/0064-heisenberg-cut}.

\bibitem{Block95} N.~Block, \emph{On a confusion about a function of
consciousness}, Behav. Brain Sci. \textbf{18}, 227 (1995).
\bibitem{Chalmers95} D.~J.~Chalmers, \emph{Facing up to the problem of
consciousness}, J. Conscious. Stud. \textbf{2}, 200 (1995).
\bibitem{Chalmers96} D.~J.~Chalmers, \emph{The Conscious Mind} (OUP, 1996).
\bibitem{DehaeneNaccache01} S.~Dehaene and L.~Naccache, \emph{Towards a
cognitive neuroscience of consciousness}, Cognition \textbf{79}, 1 (2001).
\bibitem{Dehaene14} S.~Dehaene, \emph{Consciousness and the Brain} (Viking,
2014).
\bibitem{Baars88} B.~J.~Baars, \emph{A Cognitive Theory of Consciousness}
(CUP, 1988).
\bibitem{Friston10} K.~Friston, \emph{The free-energy principle: a unified
brain theory?}, Nat. Rev. Neurosci. \textbf{11}, 127 (2010).
\bibitem{ParrFriston19} T.~Parr and K.~Friston, \emph{Generalised free energy
and active inference}, Biol. Cybern. \textbf{113}, 495 (2019).
\bibitem{Marr82} D.~Marr, \emph{Vision} (W.~H.~Freeman, 1982).
\bibitem{Simon55} H.~A.~Simon, \emph{A behavioral model of rational choice},
Q. J. Econ. \textbf{69}, 99 (1955).
\bibitem{LiederGriffiths20} F.~Lieder and T.~L.~Griffiths,
\emph{Resource-rational analysis}, Behav. Brain Sci. \textbf{43}, e1 (2020).
\bibitem{Laughlin98} S.~B.~Laughlin et al., \emph{The metabolic cost of neural
information}, Nat. Neurosci. \textbf{1}, 36 (1998).
\bibitem{AttwellLaughlin01} D.~Attwell and S.~B.~Laughlin, \emph{An energy
budget for signaling in the grey matter of the brain}, J. Cereb. Blood Flow
Metab. \textbf{21}, 1133 (2001).
\bibitem{Tononi16} G.~Tononi, M.~Boly, M.~Massimini, and C.~Koch,
\emph{Integrated information theory}, Nat. Rev. Neurosci. \textbf{17}, 450
(2016).

\end{thebibliography}

\end{document}
